\section{REST API GÜVENLİĞİ VE TEMEL KAVRAMLAR}
\label{bolum2}

\subsection{REST API Mimarisi}
\label{rest-api-mimarisi}

REST (Representational State Transfer), istemci ve sunucu arasındaki veri iletişimini HTTP protokolü üzerinden gerçekleştiren yazılım mimarisidir. REST tabanlı sistemlerde veri alışverişi genellikle JSON (JavaScript Object Notation) formatı üzerinden sağlanmaktadır. GET, POST, PUT ve DELETE gibi HTTP metodları kullanılarak istemci ile sunucu arasında kaynak odaklı veri aktarımı gerçekleştirilmektedir.

REST API yapısının temel ilkeleri şunlardır: durumsuzluk (statelessness), istemci-sunucu ayrımı, önbelleğe alınabilirlik (cacheability), katmanlı sistem (layered system) ve tek tip arayüz (uniform interface). Bu ilkeler, ölçeklenebilir ve sürdürülebilir sistem tasarımını mümkün kılmaktadır.

\subsection{Kimlik Doğrulama ve Yetkilendirme}
\label{kimlik-do-rulama-ve-yetkilendirme}

Kimlik doğrulama (authentication), sisteme giriş yapan kullanıcının gerçekten belirtilen kullanıcı olup olmadığının kontrol edilmesini ifade etmektedir. Yetkilendirme (authorization) ise doğrulanan kullanıcının hangi işlemleri gerçekleştirebileceğinin belirlenmesini sağlamaktadır. Modern API sistemlerinde bu işlemler çoğunlukla JWT (JSON Web Token) tabanlı yapılar ile gerçekleştirilmektedir.

JWT, üç bölümden oluşan imzalı bir token yapısıdır: header (algoritma bilgisi), payload (kullanıcı verileri) ve signature (imza). Kullanıcı başarıyla giriş yaptığında sunucu tarafından bir JWT üretilir ve bu token istemci tarafında saklanarak her istekte Authorization başlığı ile sunucuya iletilir. Sunucu, token imzasını doğrulayarak kullanıcının kimliğini ve yetkilerini belirler.

\subsection{OWASP API Security Top 10}
\label{owasp-top10}

OWASP API Security Top 10, REST API sistemlerinde en yaygın görülen güvenlik açıklarını tanımlayan güvenlik standardıdır. Bu proje kapsamında ele alınan başlıca zafiyetler şunlardır:

\begin{itemize}[leftmargin=*]
    \item Broken Object Level Authorization (BOLA/IDOR): Kullanıcıların başka kullanıcılara ait verilere erişebilmesi.
    \item Broken Authentication: Kimlik doğrulama mekanizmalarının yetersiz veya hatalı uygulanması.
    \item Broken Function Level Authorization: Normal kullanıcıların yönetici işlevlerine erişebilmesi.
    \item Mass Assignment: İstemci tarafından gönderilen verilerin filtrelenmeden doğrudan işlenmesi.
    \item Security Misconfiguration: Güvenlik ayarlarının yanlış veya eksik yapılandırılması.
\end{itemize}

\subsection{Savunma Mekanizmaları}
\label{savunma-mekanizmalar}

Modern API güvenliği yaklaşımlarında rate limiting, input validation ve Web Application Firewall (WAF) sistemleri önemli savunma mekanizmaları arasında yer almaktadır. Rate limiting mekanizması, istemcilerin belirli süre içerisinde gerçekleştirebileceği istek sayısını sınırlandırırken; WAF sistemleri, zararlı istekleri filtreleyerek saldırı girişimlerini engellemeyi amaçlamaktadır. Güvenli parola hashleme ve katı giriş doğrulama kontrolleri de saldırı yüzeyini önemli ölçüde azaltan temel güvenlik önlemleri arasında sayılabilir.


\clearpage
